% Copy this file to start a new recipe: cp recipes/TEMPLATE.tex recipes/my-recipe.tex
%
% The grid is a spreadsheet: ingredients go down the left column,
% prep stages (melt / mix / fold / bake, whatever your recipe needs)
% go across the top, and merged cells show what combines with what.
%
% How merges work (tabularray's `tblr`):
%   \SetCell[r=3]{c,m,font=\bfseries} mix   starts a cell that spans
%   3 rows *down* in that column. Every row after the first that the
%   merge covers must have NO entry at all for that column -- just
%   leave it out between the &'s, e.g. `ingredient & & &` for a row
%   whose 2nd/3rd columns are swallowed by merges above.
%   \SetCell[c=6]{...} spans *across* 6 columns instead (used for the
%   pan-prep / oven-preheat banner rows below).
%
% Add or remove columns by editing `colspec` and every row to match;
% add or remove ingredient rows freely -- just keep merge row-counts
% (the r=N) in sync with how many rows they actually cover.

\documentclass[10pt]{article}
\usepackage{recipe-matrix} % put this file's dir on TEXINPUTS, or `make` handles it
\pagestyle{empty}

\begin{document}

\recipetitle{Recipe Name}
\recipemeta{Yield \quad\textbullet\quad Active time}

\begin{tblr}{
  colspec = {X[2.6,l,m] Q[c,m,1.2cm] Q[c,m,1.2cm] Q[c,m,1.2cm] Q[c,m,1.5cm]},
  width = \linewidth,
  hlines = {fg=recipeaccent, wd=0.9pt},
  vlines = {fg=recipeaccent, wd=0.9pt},
  rowsep = 3pt,
  colsep = 7pt,
}
\SetCell[c=5]{c, font=\bfseries} Prep the pan / preheat the oven \\
Ingredient one & \SetCell{font=\bfseries} melt & & & \SetCell[r=4]{c,m,font=\bfseries} bake\newline 350\textdegree F\newline XX min \\
Ingredient two & & \SetCell[r=2]{c,m,font=\bfseries} mix & & \\
Ingredient three & & & & \\
Ingredient four & & & \SetCell{font=\bfseries} fold in & \\
\end{tblr}

\end{document}
